% Section 10.7, from lectures/L10/appendix/b_exercises.md.

\section{Uppgifter}
\label{c10:sec:uppgifter}

\begin{aside}
\textbf{Så arbetar du med uppgifterna.} Del~1 och del~2 räknas under lektionen: först på egen
hand, några minuter per uppgift, sedan gemensamt. Del~1 repeterar trigonometrin från \kapref{9},
och del~2 är kapitlets nya stoff. Del~3 är extrauppgifter att träna vidare på efter lektionen.

\facitnotis{L10}
\end{aside}

% ------------------------------------------------------------------------------------------
\uppgiftsdel{Del 1}{Repetitionsuppgifter}

\uppgift{1.1}{Vinklar till radianer}
Omvandla följande vinklar till radianer.
\begin{delar}(4)
  \task $15^{\circ}$
  \task $-75^{\circ}$
\end{delar}

\uppgift{1.2}{Vinklar till grader}
Omvandla följande vinklar till grader.
\begin{delar}(4)
  \task $-\pi/6$
  \task $3\pi/2$
\end{delar}

\uppgift{1.3}{Bestämning av en växelspännings egenskaper}
En växelspänning visas i figuren nedan.

\bild[0.75\linewidth]{lectures/L10/appendix/images/1.3_sine_wave.png}

Bestäm spänningens ekvation på formen $u(t) = \abs{U}\sin(\omega t + \delta)$.

\uppgift{1.4}{Ekvation samt graf för en växelspänning}
En växelspänning har amplituden $4\,\text{V}$, frekvensen $50\,\text{Hz}$ och fasen $-30^{\circ}$.
Bestäm växelspänningens ekvation $u(t)$, med fasen i radianer, och rita sinuskurvan över en
period $T$.

\uppgift{1.5}{Beräkning av en växelströms fas}
Ekvationen för en given växelspänning är
\[
  u(t) = 6\sin(80\pi t + \delta).
\]
Vid tiden $t = 15\,\text{ms}$ gäller att $u(t) = 3\,\text{V}$. Beräkna fasen $\delta$.

% ------------------------------------------------------------------------------------------
\uppgiftsdel{Del 2}{Nytt stoff}

\uppgift{2.1}{Logaritmiska ekvationer}
Lös följande logaritmiska ekvationer.
\begin{delar}(3)
  \task $5^x = 125$
  \task $2^{x-3} = 128$
  \task $e^{x+1} = 120$
\end{delar}

\uppgift{2.2}{Halveringstid för laddning i ett batteri}
Laddningen i ett batteri kan beskrivas med motsvarande spänning enligt funktionen
\[
  u(t) = U_0 \cdot a^t,
\]
där
\begin{itemize}
  \item $u(t)$ är den nuvarande laddningen,
  \item $U_0$ är den ursprungliga laddningen,
  \item $a$ är förändringsfaktorn,
  \item $t$ är antalet passerade timmar.
\end{itemize}
Ett batteri tappar halva sin laddning på $40$ timmar. Beräkna efter hur lång tid endast $10\,\%$
av laddningen återstår.

\uppgift{2.3}{Linjär förstärkning}
Två signaler har spänningsnivåerna $L_1 = 15\,\text{dB}$ och $L_2 = 49\,\text{dB}$. Beräkna den
linjära spänningsförstärkningen $G_{\text{lin}}$ mellan dessa två nivåer. Använd formeln
\[
  G_{\text{dB}} = 20\log_{10}(G_{\text{lin}})
\]
och tänk på att $G_{\text{dB}} = L_2 - L_1$.

\uppgift{2.4}{Beräkning av effektivvärde i dBV}
Sambandet mellan en spännings effektivvärde och motsvarande värde i dBV (decibel volt) är
\[
  U_{\text{dBV}} = 20\log_{10}\left(\frac{U_{\text{RMS}}}{1\,\text{V}}\right),
\]
där $U_{\text{RMS}}$ är spänningens effektivvärde i V och $U_{\text{dBV}}$ spänningen i dBV.
En sinusspänning har nivån $17{,}0\,\text{dBV}$. Bestäm amplituden i V.

% ------------------------------------------------------------------------------------------
\uppgiftsdel{Del 3}{Extrauppgifter}
Uppgifterna nedan är extra träning och görs med fördel på egen hand efter lektionen. De flesta går
att räkna i huvudet eller med papper och penna.

\uppgift{3.1}{Logaritmer utan miniräknare}
Beräkna.
\begin{delar}(6)
  \task $\log 100$
  \task $\log 1000$
  \task $\log 1$
  \task $\log 0{,}01$
  \task $\ln e$
  \task $\ln 1$
\end{delar}

\uppgift{3.2}{Logaritmlagarna}
Skriv som en enda logaritm och förenkla så långt som möjligt.
\begin{delar}(3)
  \task $\log 3 + \log 4$
  \task $\log 20 - \log 4$
  \task $2\log 5$
  \task $\log 2 + \log 5$
  \task $3\log 2 - \log 4$
\end{delar}

\uppgift{3.3}{Exponentialekvationer}
Lös ekvationerna.
\begin{delar}(5)
  \task $2^x = 32$
  \task $10^x = 0{,}001$
  \task $5^x = 100$
  \task $e^x = 20$
  \task $4^{x+1} = 64$
\end{delar}

\uppgift{3.4}{Logaritmiska ekvationer}
Lös ekvationerna.
\begin{delar}(3)
  \task $\log x = 2$
  \task $\ln x = 0$
  \task $\log(x + 5) = 1$
  \task $2\log x = 4$
  \task $\ln(2x) = 1$
\end{delar}

\uppgift{3.5}{Spänningsförstärkning i dB}
Använd
\[
  G_{\text{dB}} = 20\log_{10}\left(\frac{U_{\text{ut}}}{U_{\text{in}}}\right).
\]
\begin{parts}
  \item $U_{\text{in}} = 10\,\text{mV}$ och $U_{\text{ut}} = 1\,\text{V}$. Beräkna
    $G_{\text{dB}}$.
  \item $U_{\text{in}} = 2\,\text{V}$ och $U_{\text{ut}} = 2\,\text{V}$. Beräkna $G_{\text{dB}}$.
  \item $U_{\text{in}} = 5\,\text{V}$ och $U_{\text{ut}} = 0{,}5\,\text{V}$. Beräkna
    $G_{\text{dB}}$.
  \item Vad betyder ett negativt värde på $G_{\text{dB}}$?
\end{parts}

\uppgift{3.6}{Från dB till linjär förstärkning}
Använd
\[
  G_{\text{lin}} = 10^{G_{\text{dB}}/20}.
\]
\begin{delar}(4)
  \task $G_{\text{dB}} = 20\,\text{dB}$
  \task $G_{\text{dB}} = 60\,\text{dB}$
  \task $G_{\text{dB}} = 6\,\text{dB}$
  \task $G_{\text{dB}} = -3\,\text{dB}$
  \task* En förstärkare med $G_{\text{dB}} = 26\,\text{dB}$ matas med
    $U_{\text{in}} = 50\,\text{mV}$. Beräkna $U_{\text{ut}}$.
\end{delar}

\uppgift{3.7}{Effektförstärkning i dB}
Använd
\[
  G_{\text{dB}} = 10\log_{10}\left(\frac{P_{\text{ut}}}{P_{\text{in}}}\right).
\]
\begin{parts}
  \item $P_{\text{in}} = 1\,\text{W}$ och $P_{\text{ut}} = 100\,\text{W}$. Beräkna
    $G_{\text{dB}}$.
  \item $P_{\text{in}} = 20\,\text{W}$ och $P_{\text{ut}} = 10\,\text{W}$. Beräkna
    $G_{\text{dB}}$.
  \item Varför används faktorn $10$ för effekt men $20$ för spänning?
  \item Hur många dB motsvarar en fördubbling av effekten? Och en fördubbling av spänningen?
\end{parts}

\uppgift{3.8}{Kaskadkopplade steg}
En signalkedja består av två förstärkarsteg med $G_1 = 12\,\text{dB}$ och $G_2 = 18\,\text{dB}$,
samt en kabel som dämpar $4\,\text{dB}$.
\begin{parts}
  \item Beräkna den totala förstärkningen i dB.
  \item Beräkna den totala linjära förstärkningen.
  \item Beräkna varje delsteg linjärt och kontrollera att produkten stämmer med b).
  \item Varför är dB praktiskt att räkna med vid kaskadkoppling?
\end{parts}

\uppgift{3.9}{Nivå i dBV}
Använd
\[
  U_{\text{dBV}} = 20\log_{10}\left(\frac{U_{\text{RMS}}}{1\,\text{V}}\right).
\]
\begin{parts}
  \item $U_{\text{RMS}} = 1\,\text{V}$. Beräkna nivån i dBV.
  \item $U_{\text{RMS}} = 10\,\text{V}$. Beräkna nivån i dBV.
  \item $U_{\text{RMS}} = 0{,}1\,\text{V}$. Beräkna nivån i dBV.
  \item En signal ligger på $12\,\text{dBV}$. Beräkna $U_{\text{RMS}}$.
  \item Beräkna amplituden $\abs{U}$ för signalen i d).
\end{parts}

\uppgift{3.10}{Halverings- och fördubblingstid}
En storhet beskrivs av $f(t) = C \cdot a^t$.
\begin{parts}
  \item Ett batteri tappar $2\,\%$ av sin laddning per timme. Bestäm förändringsfaktorn $a$.
  \item Efter hur många timmar återstår hälften av laddningen?
  \item En mätdatamängd växer med $10\,\%$ per dygn. Efter hur många dygn har den fördubblats?
  \item Visa allmänt att fördubblingstiden ges av $t = \dfrac{\log 2}{\log a}$.
\end{parts}

\uppgift{3.11}{Hitta felet}
Varje rad innehåller ett vanligt fel. Förklara felet och ange det korrekta svaret.
\begin{parts}
  \item $\log(a + b) = \log a + \log b$
  \item $\dfrac{\log 8}{\log 2} = \log 4$
  \item $\log(a^n) = (\log a)^n$
  \item $G_{\text{dB}} = 0$ betyder att signalen försvinner helt.
  \item $10^{\log x} = 10x$
  \item $\ln 0 = 1$
\end{parts}
